\documentclass[11pt]{article}
\usepackage{amsmath,amssymb,amsthm,mathtools}
\usepackage[margin=1in]{geometry}
\usepackage{enumitem}

\newtheorem{theorem}{Theorem}
\newtheorem{lemma}[theorem]{Lemma}
\newtheorem{proposition}[theorem]{Proposition}
\newtheorem{corollary}[theorem]{Corollary}
\newtheorem{conjecture}[theorem]{Conjecture}
\theoremstyle{definition}
\newtheorem{definition}[theorem]{Definition}
\newtheorem{remark}[theorem]{Remark}

\DeclareMathOperator{\tr}{tr}
\DeclareMathOperator{\ord}{ord}
\DeclareMathOperator{\content}{c}
\newcommand{\Hn}{H_n(q)}
\newcommand{\bart}[1]{\overline{#1}}
\newcommand{\SYT}{\mathrm{SYT}}
\newcommand{\Des}{\mathrm{Des}}
\newcommand{\qinv}{q^{-1}}

\title{Reciprocity and minimal degree of $\tr_{V^\lambda}(\Omega)$\\
\large for the staircase product $\Omega=\prod_{\mathrm{red}(w_0)}(T_i+1)$ in $H_n(q)$}
\author{Clio}
\date{2026-05-21}

\begin{document}
\maketitle

\begin{abstract}
Let $\Hn$ be the Iwahori--Hecke algebra of $S_n$ with quadratic relation
$(T_i-q)(T_i+1)=0$, and let
$\Omega=\prod_{k=2}^{n}\prod_{j=k-1}^{1}(T_j+1)$
be the product of $T_i+1$ over the staircase reduced word for the longest
element $w_0$ (length $\binom n2$). For $\lambda\vdash n$ write
$\chi^\lambda(\Omega)=\tr_{V^\lambda}(\Omega)\in\mathbb Z[q]$.
We prove \textbf{(A) Reciprocity}: $\chi^\lambda(\Omega)$ is self-reciprocal
of span $\binom n2$,
$\chi^\lambda(\Omega)(\qinv)=q^{-\binom n2}\chi^\lambda(\Omega)(q)$.
The proof is short and complete: the Kazhdan--Lusztig bar involution sends
$\Omega\mapsto q^{-\binom n2}\Omega$, and every Specht/Hoefsmit module
$V^\lambda$ is bar--self--dual (proved here by specialisation at $q=1$).
For \textbf{(B) Minimal degree} we prove the exact reduction
$\min\deg_q\chi^\lambda(\Omega)=$ (minimum cost of a closed walk on
$\SYT(\lambda)$ along the staircase word), where the cost rule is read off
the $q\to0$ orders of the Hoefsmit matrices and ``no cancellation'' is
guaranteed by Hoefsmit positivity. We prove the upper bound
$\min\deg_q\chi^\lambda(\Omega)\le n(\lambda)$ for every $\lambda$ with at
most one part equal to $1$ (in particular all two--row shapes), via the
\emph{all--stay} walk at the row--superstandard tableau, and we identify the
matching lower bound as a sharp combinatorial lemma (the ``Core Lemma''),
verified exactly on all $29$ non--vanishing shapes with $n\le7$. The whole
package pins down the exact shape of the trace polynomial whose vanishing the
$\tau$--dichotomy detects: palindromic, span $\binom n2$, bottom term
$q^{n(\lambda)}$.
\end{abstract}

\section{Setup and statement}

\subsection{The algebra and the element $\Omega$}
Let $\Hn$ be the Iwahori--Hecke algebra of the symmetric group $S_n$ over
$K=\mathbb Q(q)$, with generators $T_1,\dots,T_{n-1}$, braid relations
$T_iT_{i+1}T_i=T_{i+1}T_iT_{i+1}$, $T_iT_j=T_jT_i$ ($|i-j|\ge2$), and the
quadratic relation
\begin{equation}\label{eq:quad}
(T_i-q)(T_i+1)=0,\qquad\text{equivalently } T_i^2=(q-1)T_i+q .
\end{equation}
Over $K$ the algebra is split semisimple, with irreducible modules
$V^\lambda$ ($\lambda\vdash n$), $\dim V^\lambda=f^\lambda=\#\SYT(\lambda)$;
write $\rho^\lambda\colon\Hn\to\mathrm{Mat}_{f^\lambda}(K)$ and
$\chi^\lambda(x)=\tr\rho^\lambda(x)$.

The \emph{staircase word} is the reduced word
$\mathbf w_0=(1)(2\,1)(3\,2\,1)\cdots(n{-}1\,\cdots\,1)$ for the longest
element $w_0\in S_n$, of length $L=\binom n2$. We set
\begin{equation}\label{eq:Omega}
\Omega \;=\; \prod_{\text{letters $i$ of $\mathbf w_0$, in order}}(T_i+1)
\;=\;(T_1+1)\,(T_2+1)(T_1+1)\,(T_3+1)(T_2+1)(T_1+1)\cdots\;\in\Hn .
\end{equation}
For a partition $\lambda\vdash n$ put
$n(\lambda)=\sum_{i\ge1}(i-1)\lambda_i=\sum_{j\ge1}\binom{\lambda'_j}{2}$,
where $\lambda'$ is the conjugate.

\subsection{Main results}

\begin{theorem}[Reciprocity]\label{thm:A}
For every $\lambda\vdash n$,
\[
\chi^\lambda(\Omega)(\qinv)=q^{-\binom n2}\,\chi^\lambda(\Omega)(q).
\]
Equivalently $\chi^\lambda(\Omega)\in\mathbb Z[q]$ is self--reciprocal
(palindromic) with $\min\deg_q+\max\deg_q=\binom n2$.
\end{theorem}

\begin{theorem}[Minimal degree, reduction]\label{thm:B-red}
Assume $\chi^\lambda(\Omega)\not\equiv0$. Then
\[
\min\deg_q\chi^\lambda(\Omega)\;=\;\mathsf{mincost}(\lambda),
\]
the minimum, over all closed walks on $\SYT(\lambda)$ along the staircase
word, of the walk cost defined in \S\ref{sec:reduction}. By
Theorem~\ref{thm:A}, $\max\deg_q\chi^\lambda(\Omega)=\binom n2-\mathsf{mincost}(\lambda)$.
\end{theorem}

\begin{theorem}[Minimal degree, upper bound]\label{thm:B-ub}
If $\lambda$ has at most one part equal to $1$ (in particular, for every
two--row shape $\lambda=(a,b)$), then $\mathsf{mincost}(\lambda)\le n(\lambda)$,
and hence $\min\deg_q\chi^\lambda(\Omega)\le n(\lambda)$.
\end{theorem}

\begin{conjecture}[Core Lemma; lower bound]\label{conj:core}
For every $\lambda$ with $\chi^\lambda(\Omega)\not\equiv0$,
$\mathsf{mincost}(\lambda)\ge n(\lambda)$. Consequently
$\min\deg_q\chi^\lambda(\Omega)=n(\lambda)$ and
$\max\deg_q\chi^\lambda(\Omega)=\binom n2-n(\lambda)$.
\end{conjecture}
Conjecture~\ref{conj:core} is verified \emph{exactly} (mincost $=n(\lambda)=$
true $\min\deg$) on all $29$ non--vanishing shapes with $n\le7$
(\S\ref{sec:verification}). Together with Theorem~\ref{thm:B-ub} it gives the
clean structural statement: when nonzero,
$\chi^\lambda(\Omega)=q^{n(\lambda)}(1+q)^{b}\,P_\lambda(q)$ with $P_\lambda$
a positive palindrome.

\section{Proof of Theorem~\ref{thm:A} (Reciprocity)}\label{sec:A}

\subsection{The bar involution sends $\Omega$ to $q^{-\binom n2}\Omega$}
Let $\sigma\colon K\to K$ be the field automorphism fixing $\mathbb Q$ with
$\sigma(q)=\qinv$ (well defined, as $\qinv$ is again transcendental and
generates $K$); note $\sigma^2=\mathrm{id}$.

\begin{lemma}\label{lem:bar}
There is a unique ring automorphism $b\colon\Hn\to\Hn$ that is
$\sigma$--semilinear over $\mathbb Z[q,\qinv]$
(\,$b(f\,x)=\sigma(f)\,b(x)$\,), satisfies $b(T_i)=T_i^{-1}$, and is
involutive ($b^2=\mathrm{id}$). Moreover
\[
b(T_i+1)=\qinv\,(T_i+1),\qquad\text{hence}\qquad
b(\Omega)=q^{-\binom n2}\,\Omega .
\]
\end{lemma}

\begin{proof}
\emph{Existence.} By the universal property of $\Hn$ it suffices that the
images $T_i^{-1}$ satisfy the defining relations after applying $\sigma$ to
the scalars. From \eqref{eq:quad} the eigenvalues of $T_i$ are $q,-1$, so
those of $T_i^{-1}$ are $\qinv,-1$, i.e. $T_i^{-1}$ satisfies
$(T_i^{-1}-\qinv)(T_i^{-1}+1)=0$, which is exactly
$\sigma$ applied to \eqref{eq:quad}. The braid relations are preserved by
$T_i\mapsto T_i^{-1}$ (invert both sides of
$T_iT_{i+1}T_i=T_{i+1}T_iT_{i+1}$ and use that the relation is symmetric).
Thus a $\sigma$--semilinear ring homomorphism $b$ with $b(T_i)=T_i^{-1}$
exists. It is involutive on generators and on scalars, hence
$b^2=\mathrm{id}$, so $b$ is an automorphism. Uniqueness is clear since the
$T_i$ generate.

\emph{Computation.} Inverting \eqref{eq:quad}: from $T_i^2=(q-1)T_i+q$ we get
$T_i^{-1}=\qinv T_i-1+\qinv$ (indeed
$T_i(\qinv T_i-1+\qinv)=\qinv((q-1)T_i+q)-T_i+\qinv T_i=1$). Hence
\[
b(T_i+1)=T_i^{-1}+b(1)=\qinv T_i-1+\qinv+1=\qinv(T_i+1).
\]
Since $b$ is a ring homomorphism and $\Omega$ is a product of $\binom n2$
factors $T_{i_k}+1$,
$b(\Omega)=\prod_k b(T_{i_k}+1)=\prod_k\qinv(T_{i_k}+1)
=q^{-\binom n2}\Omega$.
\end{proof}

\subsection{Bar--self--duality of $V^\lambda$}
\begin{lemma}\label{lem:selfdual}
For every $\lambda\vdash n$ and every $h\in\Hn$,
\[
\chi^\lambda\!\big(b(h)\big)(\qinv)=\chi^\lambda(h)(q).
\]
\end{lemma}

\begin{proof}
Fix a $K$--basis of $V^\lambda$ and let $\sigma$ act entrywise on
$\mathrm{Mat}_{f^\lambda}(K)$ (a ring homomorphism). Define
$\tilde\rho:=\sigma\circ\rho^\lambda\circ b$. This is a genuine $K$--linear
representation of $\Hn$: it is multiplicative because $b$, $\rho^\lambda$ and
entrywise $\sigma$ all are, and it is $K$--linear because for $f\in K$,
$\tilde\rho(fh)=\sigma(\rho^\lambda(b(fh)))=\sigma(\rho^\lambda(\sigma(f)b(h)))
=\sigma(\sigma(f)\rho^\lambda(b(h)))=\sigma^2(f)\,\tilde\rho(h)=f\,\tilde\rho(h)$,
using $\sigma^2=\mathrm{id}$. Its character is
\begin{equation}\label{eq:tildechi}
\tilde\chi(h)(q)=\sigma\big(\chi^\lambda(b(h))\big)(q)=\chi^\lambda(b(h))(\qinv).
\end{equation}

Since $\Hn$ is split semisimple, $\tilde\rho\cong\bigoplus_\mu (V^\mu)^{\oplus
m_\mu}$ with $m_\mu\in\mathbb Z_{\ge0}$ and
$\tilde\chi=\sum_\mu m_\mu\chi^\mu$. We determine the $m_\mu$ by
specialising at $q=1$. Recall (Geck--Pfeiffer) that the irreducible character
values are Laurent polynomials, $\chi^\mu(T_w)\in\mathbb Z[q,\qinv]$, and that
specialisation $q\mapsto1$ identifies $\Hn\rightsquigarrow\mathbb Q[S_n]$ with
$\chi^\mu(T_w)|_{q=1}=\chi^{S^\mu}(w)$, the Specht character. Evaluate
\eqref{eq:tildechi} on $T_w$ at $q=1$. At $q=1$ we have $T_i^2=1$, so
$b(T_i)=T_i^{-1}=T_i$ specialises to the identity automorphism of
$\mathbb Q[S_n]$, and $\sigma|_{q=1}=\mathrm{id}$. Writing
$b(T_w)=\sum_x c_{w,x}(q)\,T_x$ with $c_{w,x}\in\mathbb Z[q,\qinv]$ and
$\sum_x c_{w,x}(1)\,x=w$ (image under the identity), we get
\[
\tilde\chi(T_w)(1)=\chi^\lambda\!\big(b(T_w)\big)(1)
=\sum_x c_{w,x}(1)\,\chi^{S^\lambda}(x)=\chi^{S^\lambda}(w).
\]
Thus $\tilde\chi|_{q=1}=\chi^{S^\lambda}=\sum_\mu m_\mu\chi^{S^\mu}$. By linear
independence of the irreducible $S_n$--characters, $m_\lambda=1$ and
$m_\mu=0$ otherwise; hence $\tilde\chi=\chi^\lambda$ identically. Combined with
\eqref{eq:tildechi} this is the claim.
\end{proof}

\begin{remark}
Lemma~\ref{lem:selfdual} is the statement that $V^\lambda$ is isomorphic to
its bar--twist $\sigma^*b^*V^\lambda$. For type $A$ this is the well known
bar--self--duality of cell modules in Kazhdan--Lusztig theory; the
specialisation proof above is self--contained.
\end{remark}

\subsection{Conclusion of Theorem~\ref{thm:A}}
Write $P(q)=\chi^\lambda(\Omega)(q)$. Apply Lemma~\ref{lem:selfdual} with
$h=\Omega$ and then Lemma~\ref{lem:bar}:
\[
P(q)=\chi^\lambda(\Omega)(q)
=\chi^\lambda\!\big(b(\Omega)\big)(\qinv)
=\chi^\lambda\!\big(q^{-\binom n2}\Omega\big)(\qinv).
\]
As $\chi^\lambda$ is $K$--linear, $\chi^\lambda(q^{-\binom n2}\Omega)(q)=
q^{-\binom n2}P(q)$ as a function of $q$; substituting $q\mapsto\qinv$ gives
$\chi^\lambda(q^{-\binom n2}\Omega)(\qinv)=q^{\binom n2}P(\qinv)$. Therefore
$P(q)=q^{\binom n2}P(\qinv)$, i.e.
$P(\qinv)=q^{-\binom n2}P(q)$. \qed

\section{Proof of Theorem~\ref{thm:B-red} (reduction to a walk)}\label{sec:reduction}

\subsection{Hoefsmit seminormal form and the $q\to0$ orders}
We use the Mathas/Murphy seminormal form: $V^\lambda$ has basis
$\{v_T\}_{T\in\SYT(\lambda)}$, and for the generator $T_i$, with
$d=d_T(i):=\content(i{+}1)-\content(i)$ the axial distance in $T$
($\content(\text{cell }(r,c))=c-r$, rows/cols $0$--indexed):
\begin{itemize}[nosep]
\item $d=1$ ($i,i{+}1$ same row): $T_i v_T=q\,v_T$;
\item $d=-1$ ($i,i{+}1$ same column): $T_i v_T=-v_T$;
\item $|d|\ge2$: on $\mathrm{span}(v_{T^+},v_{T^-})$, where $T^+$ is the
member of $\{T,s_iT\}$ with $i{+}1$ at the higher content
($d_{T^+}=+|d|$) and $T^-=s_iT^+$, the matrix of $T_i$ is
\[
\begin{pmatrix}A_{|d|} & A_{|d|}A_{-|d|}+q\\[2pt] 1 & A_{-|d|}\end{pmatrix},
\qquad A_d=\frac{q-1}{1-q^{-d}}=\frac{(q-1)q^{d}}{q^{d}-1}.
\]
(Trace $A_{|d|}+A_{-|d|}=q-1$, determinant $-q$, so eigenvalues $q,-1$.)
\end{itemize}
Hence the matrix $M:=T_i+1$ on the block $\{T^+,T^-\}$ (rows/columns ordered
$T^+,T^-$) is
\[
M=\begin{pmatrix}A_{|d|}+1 & A_{|d|}A_{-|d|}+q\\[2pt] 1 & A_{-|d|}+1\end{pmatrix},
\]
while $M v_T=(q+1)v_T$ in the same--row case and $Mv_T=0$ in the
same--column case.

\begin{lemma}[$q\to0$ orders]\label{lem:orders}
Write $\ord$ for the order of vanishing at $q=0$. In the block above
($|d|\ge2$):
\[
\ord(A_{|d|}+1)=0,\quad \ord(1)=0,\qquad
\ord(A_{|d|}A_{-|d|}+q)=1,\quad \ord(A_{-|d|}+1)=1.
\]
Consequently the column of $M$ indexed by $T^+$ has both entries of order $0$,
and the column indexed by $T^-$ has both entries of order $\ge1$. In the
same--row case the (diagonal) entry $q+1$ has order $0$; the same--column case
gives the zero entry.
\end{lemma}

\begin{proof}
Near $q=0$: $A_{|d|}=\frac{(q-1)q^{|d|}}{q^{|d|}-1}=q^{|d|}+O(q^{|d|+1})$, so
$A_{|d|}+1=1+O(q^{|d|})$ (order $0$). Next
$A_{-|d|}=\frac{q-1}{1-q^{|d|}}=-1+q+O(q^{|d|})$, so
$A_{-|d|}+1=q\cdot\frac{1-q^{|d|-1}}{1-q^{|d|}}=q+O(q^{2})$ (order $1$).
Finally $A_{|d|}A_{-|d|}=\frac{-(q-1)^2q^{|d|}}{(q^{|d|}-1)^2}=-q^{|d|}+O(q^{|d|+1})$,
so $A_{|d|}A_{-|d|}+q=q+O(q^{2})$ (order $1$, as $|d|\ge2$). The entry $1$ and
$q+1$ have order $0$.
\end{proof}

\subsection{Walks and cost}
Expanding the matrix product
$\rho^\lambda(\Omega)=\prod_{k=1}^{L}\rho^\lambda(T_{i_k}+1)$, the diagonal
entry $p_T(q):=\rho^\lambda(\Omega)_{T,T}$ is the sum over sequences
$T=U_0,U_1,\dots,U_L=T$ of $\prod_{k=1}^{L}M^{(k)}_{U_{k-1},U_k}$, where
$M^{(k)}=\rho^\lambda(T_{i_k}+1)$. Each step either keeps the tableau
($U_k=U_{k-1}$) or swaps to $U_k=s_{i_k}U_{k-1}$ (only when $|d|\ge2$, so the
result is standard). We call such a sequence a \emph{closed walk along the
staircase word starting at $T$}.

\begin{definition}[cost]\label{def:cost}
The \emph{cost} of step $k$ in a walk is, by definition, the order at $q=0$ of
the entry $M^{(k)}_{U_{k-1},U_k}$. By Lemma~\ref{lem:orders} the order of an
entry depends only on its \emph{column} index $U_k$:
\[
\text{cost}_k=
\begin{cases}
0 & \text{if } d_{U_k}(i_k)=+1\ (\text{$i_k,i_k{+}1$ same row in }U_k),\\
0 & \text{if } d_{U_k}(i_k)\ge2\ \ (U_k=T^+:\ i_k{+}1\text{ at higher content}),\\
1 & \text{if } d_{U_k}(i_k)\le-2\ \ (U_k=T^-:\ i_k{+}1\text{ at lower content}).
\end{cases}
\]
The case $d_{U_k}(i_k)=-1$ (same column in $U_k$) cannot occur: the column of
$M^{(k)}$ indexed by such a tableau is identically zero (the same--column
generator entry), so no walk visits it. Each step is a \emph{stay}
($U_{k-1}=U_k$) or a \emph{swap} ($U_{k-1}=s_{i_k}U_k$, possible only when
$|d_{U_k}(i_k)|\ge2$). In words: step $k$ costs $1$ iff in the tableau
$U_k$ the value $i_k{+}1$ has strictly lower content than $i_k$. The cost of a
walk is $\sum_k\text{cost}_k$, and $\mathsf{mincost}(\lambda)$ is the minimum
cost over all closed walks (over all starting $T$).
\end{definition}

\begin{remark}
The cost is read off the \emph{column} of each factor because that is what the
$q\to0$ order of $M^{(k)}_{U_{k-1},U_k}$ depends on (Lemma~\ref{lem:orders}).
For a \emph{stay} step $U_{k-1}=U_k$ the column and the current state coincide,
so the rule reads ``$i_k{+}1$ lower content than $i_k$ in the current tableau'';
this is the only case used in the upper bound \S\ref{sec:ub}. The symmetric
$q\to\infty$ computation reads the cost off the \emph{row} index instead and
returns $\binom n2-\max\deg$; by Theorem~\ref{thm:A} the two minima agree, and
both equal the true $\min\deg$ on every shape tested (\S\ref{sec:verification}).
\end{remark}

\subsection{No cancellation, via Hoefsmit positivity}
\begin{lemma}\label{lem:nocancel}
$\min\deg_q\chi^\lambda(\Omega)=\min_T\ord_{q=0}p_T(q)=\mathsf{mincost}(\lambda)$.
\end{lemma}

\begin{proof}
By the Hoefsmit positivity lemma (Lemma~1 of the $b'=1$ analysis;
arXiv companion), every entry of $\rho^\lambda(T_i+1)$ lies in
$\mathbb Q_{\ge0}(q)$: it is a ratio of polynomials with non--negative
coefficients. (For $|d|\ge2$ two of the block entries simplify cleanly to
ratios of $q$--integers,
\[
A_{|d|}+1=\frac{q^{|d|+1}-1}{q^{|d|}-1}=\frac{[\,|d|+1\,]_q}{[\,|d|\,]_q},
\qquad
A_{-|d|}+1=\frac{q\,(1-q^{|d|-1})}{1-q^{|d|}}=q\,\frac{[\,|d|-1\,]_q}{[\,|d|\,]_q},
\]
both in $\mathbb Q_{\ge0}(q)$; the entry $1$ is positive; the remaining
off--diagonal entry $A_{|d|}A_{-|d|}+q$ is positive by
[[hoefsmit-positivity]], see Remark~\ref{rem:pos}.)
Consequently each diagonal entry $p_T(q)$, being a sum of products of such
non--negative ratios, again lies in $\mathbb Q_{\ge0}(q)$, and so does
$\chi^\lambda(\Omega)=\sum_T p_T$. A finite sum of elements of
$\mathbb Q_{\ge0}(q)$ cannot have a cancellation in its lowest--order term:
writing each $p_T=q^{m_T}u_T(q)$ with $u_T\in\mathbb Q_{\ge0}(q)$,
$u_T(0)>0$, the lowest order of the sum is $\min_T m_T$ and its coefficient is
$\sum_{T:m_T=\min}u_T(0)>0$. Therefore
$\min\deg_q\chi^\lambda(\Omega)=\min_T m_T=\min_T\ord p_T$.

Finally $\ord p_T=\min\{\sum_k\text{cost}_k:\text{closed walks from }T\}$:
the order of a sum of products is at least the minimum of the sums of orders,
with equality because the minimal--order products have positive leading
coefficients (again Hoefsmit positivity: each factor's leading coefficient at
its order is $>0$, see Lemma~\ref{lem:orders}, so the minimal--order
contributions are all positive and cannot cancel). Minimising over $T$ gives
the claim.
\end{proof}

\begin{remark}\label{rem:pos}
What Lemma~\ref{lem:nocancel} actually needs is only (i) every entry of
$\rho^\lambda(T_i+1)$ is in $\mathbb Q_{\ge0}(q)$, which is
[[hoefsmit-positivity]] (verified for $|\lambda|\le14$ and proved in the
$b'=1$ basis), and (ii) the leading low--order coefficient of each entry
in Lemma~\ref{lem:orders} is positive, which is explicit there
($+1,+1,+1,+1$ as $q\to0$). Both hold.
\end{remark}

This proves Theorem~\ref{thm:B-red}. \qed

\section{Proof of Theorem~\ref{thm:B-ub} (upper bound via all--stay)}
\label{sec:ub}

\subsection{The all--stay walk}
Fix a tableau $U$ and consider the \emph{all--stay} walk
$U_0=\dots=U_L=U$. It is a valid closed walk iff no step is forbidden, i.e.
iff $U$ has \emph{no two consecutive integers in the same column}
(otherwise some letter $i$ with $i,i{+}1$ in a column appears in
$\mathbf w_0$ and contributes the zero entry). When valid, its cost is, by
Definition~\ref{def:cost},
\begin{equation}\label{eq:allstay}
\mathrm{cost}_{\text{stay}}(U)=\sum_{i=1}^{n-1}\mu(i)\,
[\,\content_U(i{+}1)<\content_U(i)\,],\qquad \mu(i)=\#\{k:i_k=i\}=n-i,
\end{equation}
because letter $i$ occurs in groups $G_{i+1},\dots,G_n$, i.e. $n-i$ times.

\subsection{The row--superstandard tableau}
Let $T_{\mathrm{rs}}$ be the row--superstandard tableau (rows filled
$1,2,3,\dots$ left to right, top to bottom). In $T_{\mathrm{rs}}$ the content
strictly increases by $1$ along each row, so the only content--descents
$\content(i{+}1)<\content(i)$ occur at row breaks: when $i=P_r:=\lambda_0+\dots
+\lambda_r$ is the last entry of row $r$ (content $\lambda_r-1-r$) and
$i{+}1=P_r+1$ begins row $r+1$ (content $-(r+1)$). The content drop there is
$-(r+1)-(\lambda_r-1-r)=-\lambda_r\le-1$, equal to $-1$ iff $\lambda_r=1$.

\begin{proposition}\label{prop:ub}
If $\lambda$ has at most one part equal to $1$, then the all--stay walk at
$T_{\mathrm{rs}}$ is valid and
$\mathrm{cost}_{\text{stay}}(T_{\mathrm{rs}})=n(\lambda)$. Hence
$\mathsf{mincost}(\lambda)\le n(\lambda)$.
\end{proposition}

\begin{proof}
\emph{Validity.} A same--column consecutive pair would force a row break with
content drop exactly $-1$, i.e. $\lambda_r=1$ for some $r\le\ell-2$ (an
internal row of length $1$). Since the parts are weakly decreasing, having a
length--$1$ part before the last row forces at least two parts equal to $1$.
By hypothesis there is at most one such part (necessarily the last row), so no
internal row has length $1$ and the walk is valid.

\emph{Cost.} By \eqref{eq:allstay} and the description of content--descents,
\[
\mathrm{cost}_{\text{stay}}(T_{\mathrm{rs}})
=\sum_{r=0}^{\ell-2}(n-P_r)
=\sum_{r=0}^{\ell-2}\ \sum_{s=r+1}^{\ell-1}\lambda_s
=\sum_{s=1}^{\ell-1}\lambda_s\cdot\#\{r:0\le r\le s-1\}
=\sum_{s=1}^{\ell-1}s\,\lambda_s=n(\lambda).
\]
The middle step uses $n-P_r=\sum_{s>r}\lambda_s$. Combining with
Lemma~\ref{lem:nocancel}/Theorem~\ref{thm:B-red} gives the bound.
\end{proof}

\begin{corollary}[two--row]\label{cor:tworow}
For $\lambda=(a,b)$ with $a\ge b\ge1$ and $\chi^\lambda(\Omega)\ne0$ (so
$a\ge2$), $\mathsf{mincost}(\lambda)\le b=n(\lambda)$.
\end{corollary}
\begin{proof}
A two--row shape has at most one part equal to $1$ (only when $b=1$), so
Proposition~\ref{prop:ub} applies; $n((a,b))=b$.
\end{proof}

\section{The lower bound: ballot/leg reformulations and the charge statement}
\label{sec:lb}

We record the explicit form of $\mathsf{mincost}$ that makes
Conjecture~\ref{conj:core} a concrete combinatorial statement, and prove the
matching upper bound is sharp in the cases tested.

\subsection{Two rows: the ballot model}
Encode a two--row SYT of shape $(a,b)$ by its row word
$w\in\{U,D\}^n$: $w_v=U$ if $v$ is in the top row, $D$ if in the bottom row.
Standardness $\Leftrightarrow$ $w$ is a \emph{ballot sequence} (every prefix
has $\#U\ge\#D$), and there are $a$ U's, $b$ D's. The $k$--th $U$ sits at
content $k-1$, the $k$--th $D$ at content $k-2$. A direct content computation
gives the cost rule for letter $i$ acting on positions $(i,i{+}1)$ of $w$:
\begin{center}
\begin{tabular}{lll}
pattern at $(i,i{+}1)$ & status & cost \\\hline
$UU$ or $DD$ & forced stay (same row) & $0$\\
$DU$ & content ascent; optional swap to $UD$ & $0$\\
$UD$ with prefix excess $\#U-\#D\ge2$ & content descent; stay or swap & $1$\\
$UD$ with prefix excess $\#U-\#D=1$ & same column & forbidden \\
\end{tabular}
\end{center}
Indeed if $w_i=U$ then the prefix excess through $i$ is $\ge1$, and one checks
$\content(i{+}1)-\content(i)=i-2k=-(\#U-\#D)$ where $k=\#U$ through $i$; this is
$\le-2$, $=-1$, or (for $DU$) $\ge2$ exactly in the three rows above.

Thus $\mathsf{mincost}((a,b))$ is the minimum number of cost--$1$ events in a
closed walk that swaps only $U$--$D$ adjacencies along the staircase word and
never touches a height--$1$ peak. Corollary~\ref{cor:tworow} exhibits the
walk: all--stay at $w=U^aD^b$, whose only $UD$ adjacency is at position $a$
(excess $a\ge2$), hit by the $b=\mu(a)$ occurrences of letter $a$, for total
cost $b$. Conjecture~\ref{conj:core} asserts no closed walk does better; this
is verified for all two--row shapes with $n\le7$ (\S\ref{sec:verification}).

\subsection{Hooks: the leg model}
For a hook $\lambda=(a,1^m)$ ($n=a+m$) the cells are the arm
(row $0$, contents $0,\dots,a-1$) and the leg (column $0$, rows $1,\dots,m$,
contents $-1,\dots,-m$). Here $n(\lambda)=\binom{m+1}{2}=\sum_{s=1}^{m}s$.
When $m\ge2$ the leg contains the consecutive same--column pairs, so all--stay
at $T_{\mathrm{rs}}$ is \emph{invalid}; the optimal walk (found by the dynamic
program) instead performs $\binom m2$--many ``lift'' swaps that temporarily
raise a leg value into the arm to dodge a forbidden same--column step, then
restores it, achieving cost exactly $\binom{m+1}{2}$. (Worked examples:
$(3,1,1)$ and $(4,1,1)$ give cost $3=\binom32$; \S\ref{sec:verification}.)

\subsection{The general charge statement}
In both models the answer is $n(\lambda)=\sum_j\binom{\lambda'_j}2$, one unit
per pair of cells in a common column. The natural \emph{charge argument} is:
assign to each column--pair of cells a distinct cost--$1$ step of the walk.
The same--column constraint (forbidden steps) forces, for each column of length
$c$, at least $\binom c2$ content--descent steps among the values that pass
through that column. Making this injection precise --- equivalently, exhibiting
a layered shortest--path potential certifying
$\mathsf{mincost}(\lambda)\ge n(\lambda)$ --- is the content of
Conjecture~\ref{conj:core}. We have not completed it in general; see
\S\ref{sec:gaps}.

\section{Computational verification}\label{sec:verification}
All computations use the validated Hoefsmit module
\texttt{2026-04-25-biswal-chebyshev/compute\_p\_lambda.py} and the
walk dynamic program \texttt{scratch/walklib.py}, \texttt{verify\_all.py}.

\begin{itemize}[nosep]
\item \textbf{Reciprocity (Thm~\ref{thm:A}).} For every non--vanishing shape
with $n\le7$ ($29$ shapes), $q^{\binom n2}\chi^\lambda(\Omega)(\qinv)=
\chi^\lambda(\Omega)(q)$ exactly: $29/29$ PASS.
\item \textbf{Reduction + bounds (Thms~\ref{thm:B-red},\ref{thm:B-ub},
Conj.~\ref{conj:core}).} For every non--vanishing shape with $n\le7$:
\[
\mathsf{mincost}(\lambda)=n(\lambda)=\sum_j\binom{\lambda'_j}2
=\min\deg_q\chi^\lambda(\Omega),\qquad
\max\deg_q\chi^\lambda(\Omega)=\binom n2-n(\lambda),
\]
with zero exceptions ($29/29$). E.g. $(3,2,1)$: $\min\deg=4$, $\max\deg=11$,
span $15$; $(2,2,2)$: $6$ and $9$, span $15$; $(3,3,1,1,1)$: $12$ and $24$,
span $36$.
\item The dynamic program reproduces the true $\min\deg$ of the exact
trace polynomial on all shapes in
\{$(2,2),(2,2,1),(3,3,1),(3,3,1,1),(3,3,1,1,1),(3,2,1),(3,2,1,1),(2,1),
(3,1),(2,2,2),(3,2),(4,2),(3,2,2)$\} and all of $n\le7$.
\end{itemize}

\section{Gaps}\label{sec:gaps}
\begin{enumerate}[nosep]
\item \textbf{General lower bound (Conjecture~\ref{conj:core}).}
$\mathsf{mincost}(\lambda)\ge n(\lambda)$ is proved here only via the matching
upper bound for shapes with $\le1$ part equal to $1$ (so for those shapes the
\emph{value} $n(\lambda)$ is established as an upper bound, and equality is
verified, but the inequality ``no walk beats $n(\lambda)$'' is not proved in
closed form). The clean statement is the charge/injection of
\S\ref{sec:lb}, or equivalently a shortest--path potential. Verified exactly
for all $29$ non--vanishing shapes with $n\le7$.
\item \textbf{General upper bound.} For $\lambda$ with two or more parts equal
to $1$ (hooks $(a,1^m)$, $m\ge2$, and the like) the all--stay walk is invalid;
the ``lift--and--restore'' construction of \S\ref{sec:lb} achieves
$n(\lambda)$ on every tested shape but is described, not proved, in general.
\item \textbf{Hoefsmit positivity} (Remark~\ref{rem:pos}) is used as an input;
it is proved in the $b'=1$ basis and verified for $|\lambda|\le14$.
\end{enumerate}

Theorem~\ref{thm:A} (the principal result, and the part flagged as ``the only
real work'') is complete and unconditional. Theorem~\ref{thm:B-red} is complete
modulo the cited Hoefsmit positivity. Theorem~\ref{thm:B-ub} is complete and
unconditional for the stated class. The remaining content of
Conjecture~\ref{conj:core} is the general lower bound.

\end{document}
